\section{校正性质与损失导数}
\label{app:math}
对同一帧、同一通道，由 $z=y-\hat\delta$ 得到 $z-\bar z=y-\bar y$，表明进一步处理前的通道内空间样本方差同样保持不变。由于 $\hat\delta$ 来自含噪数据，重复拍摄时有
\begin{align}
 \operatorname{Cov}(z_p,z_q)
 &=\operatorname{Cov}(y_p,y_q)+\operatorname{Var}(\hat\delta)\nonumber\\
 &\quad-\operatorname{Cov}(y_p,\hat\delta)
       -\operatorname{Cov}(y_q,\hat\delta).
\end{align}
该协方差包含估计量的方差及其与各输入样本的相关项。式\eqref{eq:invariance}描述带符号 RAW 校正阶段的帧内空间差分，后续裁切、再量化、白平衡与去马赛克均作用于这一校正信号。

为推导式\eqref{eq:gradient}，令 $u=2s/B$、$q=(\delta/B+1)/2$，且 $|\delta|\le B$，则
\begin{align}
 \frac{\partial[B^2\BCE(u,q)]}{\partial s}
 &=2B\left[\sigma(2s/B)-q\right]\nonumber\\
 &=B\tanh(s/B)-\delta.
\end{align}
网络参数梯度还包含网络雅可比、标签平均及三项损失权重，输出导数为 $\operatorname{sech}^2(s/B)$。在网络输出值 $s$ 固定时，将边界从 $B=1$ 改为 $B=3$ 仍会微调前向映射曲线，在初始评估批次上的最大输出差异为 0.0530587 DN。表\ref{tab:metrics}中的基线结果统一在扩界后的 $B=3\DN$ 下测定。

\section{训练、扰动实验与模型选择}
\label{app:protocol}
采样区域的高宽均为目标 Bayer 裁块的三倍，对应 $1.5\times$ 视野与每轴两倍积分，合格源图像的原生分辨率满足该尺寸要求。亮度缩放使用源图像 RGB 联合中位数确定的统一乘数，保留逐帧与逐 CFA 均值之间由场景决定的关系。

共同偏移符号在每个 CFA 上独立、等概率采样并在七帧内保持不变，标签由该分量与联合经验轨迹相加得到，保留共同水平和原始时间轨迹，并以零偏移序列和恒定偏移序列提供参照条件。适配保留 195 个参数张量的 Adam 优化器状态、随机数生成器状态及冻结参数集合。预训练单帧模型对中间和最终预测头分别赋予 $0.25$ 和 $0.75$ 的损失权重，并在三个独立采样的单帧视图上取平均。

图\ref{fig:designed_regimes}在同一段 120 帧配对开发序列上区分四类偏移工况，涵盖 $\delta_{t,c}=0$、$m_c$、$m_c+r_{t,c}$ 及 $m_c+r_{t,c}+j_{t,c}$。逐 CFA 指定本底为 $m=(0.60,-0.45,0.35,-0.55)\DN$，依次对应 R、G1、G2、B 通道，联合轨迹 $r$ 来自 Z63 暗场实测状态转移模型，而 $j$ 在第 24、42、59、78、96、105 帧引入六次指定的通道跳变。各类条件共用场景几何及光子与读出噪声抽样。

\figwide{designed_offset_regimes_zh}{.98}{105k 权重在合成 Bayer 评估序列上的设计偏移条件。左列为四 CFA 注入真值，右列为 R 通道真值、同权重 T1、T7 及 T1 中心七帧均值。面板内 $A/J$ 为 108 个有效中心帧上的四通道平均值（单位 DN），竖虚线标记设计变化时刻。}{fig:designed_regimes}

评估协议在每个场景、每条实测偏移轨迹下包含 2 DN 强度的七种条件（即 empirical、native-motion、zero、common、common-plus、relative 与 relative-plus，其中 native-motion 移除共同注入而保留经验变化），并在 0.25、16、64 DN 各评价 zero、common、common-plus 三种条件，每个源图像与偏移轨迹组合共 16 组。

共同偏移实验共享场景、几何和噪声种子，比较幅度为 0.3 与 0.5 DN 的有符号共同偏移。相对偏移实验在有符号 0.4-DN 恒定偏移上加入有符号轨迹 $[0,.2,-.2,.2,-.2,.2,0]\DN$。令 $\Delta p$、$\Delta y$ 为配对预测和目标差，每个 CFA 的共同增益为 $\mean_t\Delta p/\mean_t\Delta y$，用波浪号表示沿时间中心化，则相对增益为
\begin{equation}
 g_c=\frac{\sum_t\widetilde{\Delta p}_{t,c}\widetilde{\Delta y}_{t,c}}
 {\sum_t\widetilde{\Delta y}_{t,c}^{\,2}}.
\end{equation}
报告值在各源图像与偏移轨迹组合之间平均。扰动先于量化，连续偏移根据各像素量化前的数值映射为整数像素差。

早期基线选模准则要求 $C$、$J$ 与零偏移 MAE 相对初始基线上浮不超过 0.005 DN，零偏移 P95 上浮不超过 0.01 DN，再按 $(A_{T1}+A_{T7})/2$ 排序，据此选出 46k 检查点。经过开发评估，本文将 105k 检查点固定为时域分析工作点，权衡绝对偏移精度与动态跟踪响应。在 102k、105k 和 106k 检查点下，T7 的 $A/J$ 分别为 0.098377/0.055296、0.102115/0.054778 与 0.096073/0.052858 DN。虽然 106k 取得了略低的误差数值，但其在 2 DN 中位亮度下的最小共同与相对增益分别为 0.655 与 0.496，而 105k 检查点分别为 0.745 与 0.505。图\ref{fig:operating}展示了各检查点在训练过程中的精度与稳定性演变轨迹。

\section{PTC 记录与偏移模式}
\label{app:ptc}
\enlargethispage{10pt}
\begin{inlinetable}
\centering
\caption{图\ref{fig:ptc} 的光子转移曲线参数（读出噪声 RN 单位为 DN，转换增益 $k$ 单位为 DN/电子）。}\label{tab:ptc}
\footnotesize
\renewcommand{\arraystretch}{0.82}
\setlength{\tabcolsep}{3pt}
\begin{tabular}{llrrr}
\toprule
记录 & 参数来源相机 & ISO & RN & $k$ \\
\midrule
P01 & Canon EOS R5 & 400 & 1.9270 & 1.2613 \\
P02 & Canon EOS-1D X Mark II & 800 & 4.1892 & 1.1122 \\
P03 & Nikon D4s & 100 & 1.7920 & 0.1264 \\
P04 & Canon EOS RP & 400 & 9.4564 & 0.7756 \\
P07 & Sony Alpha 850 & 100 & 5.4993 & 0.5958 \\
P09 & Pixii Max & 400 & 5.4000 & 0.6614 \\
P10 & Sony A1 & 200 & 3.3464 & 0.7931 \\
P11 & Nikon D800 & 400 & 4.0161 & 1.1888 \\
P12 & Sony Alpha 290 & 100 & 6.9718 & 0.6632 \\
P14 & Nikon D300s & 200 & 4.5118 & 0.7525 \\
P15 & Pentax Kr & 200 & 5.6799 & 0.9775 \\
P17 & Pentax K-5 IIs & 100 & 1.3530 & 0.4211 \\
P19 & Olympus OM-D E-M5 & 400 & 4.7714 & 1.2557 \\
P20 & Panasonic Lumix DMC-GH4 & 400 & 5.6856 & 1.5518 \\
\bottomrule
\end{tabular}
\end{inlinetable}

\begin{inlinetable}
\centering
\caption{P01 的偏移标签示例（DN），通道顺序同图\ref{fig:ptc}。}\label{tab:patterns}
\footnotesize
\setlength{\tabcolsep}{4pt}
\begin{tabular}{lrrrr}
\toprule
模式 & R & $G_1$ & $G_2$ & B \\
\midrule
zero & 0.0000 & 0.0000 & 0.0000 & 0.0000 \\
common\_plus & 0.7378 & 0.7378 & 0.7378 & 0.7378 \\
common\_minus & -0.7378 & -0.7378 & -0.7378 & -0.7378 \\
mixed & 0.7809 & -0.9499 & 0.6017 & -0.6155 \\
mixed\_reverse & -0.7809 & 0.9499 & -0.6017 & 0.6155 \\
small\_plus & 0.2500 & 0.2500 & 0.2500 & 0.2500 \\
small\_minus & -0.2500 & -0.2500 & -0.2500 & -0.2500 \\
large\_plus & 0.9500 & 0.9500 & 0.9500 & 0.9500 \\
large\_minus & -0.9500 & -0.9500 & -0.9500 & -0.9500 \\
unequal\_plus & 0.1500 & 0.4000 & 0.7000 & 0.9500 \\
unequal\_minus & -0.1500 & -0.4000 & -0.7000 & -0.9500 \\
three\_plus & 0.8500 & 0.6500 & 0.4500 & -0.7500 \\
three\_minus & -0.8500 & -0.6500 & -0.4500 & 0.7500 \\
\bottomrule
\end{tabular}
\end{inlinetable}

表\ref{tab:ptc}列出矩阵行标对应的噪声参数记录，当前模拟中所有记录共用 15871-DN 动态范围。表\ref{tab:patterns}展示 P01 记录的具体偏移模式。尽管不同记录的随机幅值有所差异，所有 182 个测试单元的精确标签与模型预测均完整保存在配套数据中，此处 1 DN 表示统一 14 位模拟域内的一个码值单位。

单帧模型在各测试单元上接受评测，对比相同噪声实现下的校正结果与原始输入。

\section{逐场景结果与统计方法}
\label{app:statistics}
\figwide{operating_points_zh}{.98}{八场景开发集上的 $A$--$J$ 关系。轨迹包含 70k--105k 的评价点，各标记分别指出早期 46k、核心 105k 和后续 106k 检查点。每个点的 T1 与 T7 使用相同权重。}{fig:operating}
\figwide{z63_inference}{.93}{46k 检查点在 Z63 域合成序列上的视觉对比示例。场景中位亮度为 0.25 DN，显示增益固定为 +3 EV。四通道注入偏移范数最大的有效帧展示未校正、offline T1、滑动 T7 和已知偏移校正，所有面板共用相机运动和噪声。}{fig:z63_inference_archive}
\begin{inlinetable}
\centering
\caption{105,000 次适配迭代后的逐场景指标（DN），每个场景由两条偏移轨迹的测量值取平均。}\label{tab:sources}
\footnotesize
\setlength{\tabcolsep}{3pt}
\begin{tabular}{lrrrrrr}
\toprule
场景 & $A$: T1 & $A$: T7 & $J$: T1 & $J$: T7 & $S$: T1 & $S$: T7 \\
\midrule
1 & 0.0757 & 0.0681 & 0.0744 & 0.0595 & 0.2455 & 0.2399 \\
2 & 0.1554 & 0.1536 & 0.0803 & 0.0608 & 0.1198 & 0.1039 \\
3 & 0.1395 & 0.1408 & 0.0595 & 0.0503 & 0.1096 & 0.1018 \\
4 & 0.1259 & 0.1235 & 0.0490 & 0.0490 & 0.1654 & 0.1677 \\
5 & 0.0738 & 0.0676 & 0.0978 & 0.0640 & 0.2136 & 0.1922 \\
6 & 0.0779 & 0.0701 & 0.0603 & 0.0402 & 0.1294 & 0.1209 \\
7 & 0.1523 & 0.1400 & 0.0759 & 0.0477 & 0.1804 & 0.1584 \\
8 & 0.0697 & 0.0532 & 0.0921 & 0.0669 & 0.1516 & 0.1185 \\
\bottomrule
\end{tabular}
\end{inlinetable}

\begin{inlinetable}
\centering
\caption{105,000 次适配迭代后的扰动响应增益。数值 0.25、2、16 与 64 表示源图像中位强度（DN），增益为无量纲量，相对增益采用时间中心化。}\label{tab:gains}
\footnotesize
\setlength{\tabcolsep}{4pt}
\begin{tabular}{lrrrr}
\toprule
扰动／推断方式 & R & $G_1$ & $G_2$ & B \\
\midrule
共同 0.25 / T1 & 0.6211 & 0.6059 & 0.6017 & 0.6720 \\
共同 0.25 / T7 & 0.6125 & 0.5982 & 0.5930 & 0.6685 \\
共同 2 / T1 & 0.7967 & 0.8218 & 0.8158 & 0.8311 \\
共同 2 / T7 & 0.7454 & 0.7626 & 0.7577 & 0.7476 \\
共同 16 / T1 & 0.5654 & 0.4591 & 0.4711 & 0.5585 \\
共同 16 / T7 & 0.5702 & 0.4628 & 0.4663 & 0.5723 \\
共同 64 / T1 & 0.1402 & 0.0651 & 0.0750 & 0.1975 \\
共同 64 / T7 & 0.0934 & 0.0763 & 0.0796 & 0.1621 \\
相对 2 / T1 & 0.7887 & 0.7893 & 0.7849 & 0.8501 \\
相对 2 / T7 & 0.6083 & 0.5046 & 0.5242 & 0.5850 \\
\bottomrule
\end{tabular}
\end{inlinetable}

\begin{inlinetable}
\centering
\caption{105,000 次适配迭代后的经验偏移子集误差统计，按符号与幅度区间划分。MAE 单位为 DN，W 为相对于未校正时误差增大的样本比例。各分量来自八个场景。}\label{tab:buckets}
\footnotesize
\setlength{\tabcolsep}{3pt}
\begin{tabular}{lrrrrr}
\toprule
偏移组 & $N$ & T1 MAE & T7 MAE & T1 W & T7 W \\
\midrule
negative & 262 & 0.1138 & 0.1076 & 3.44\% & 3.05\% \\
positive & 186 & 0.1018 & 0.0944 & 10.75\% & 10.75\% \\
$(0,.1]$ & 13 & 0.1107 & 0.1092 & 46.15\% & 53.85\% \\
$(.1,.3]$ & 93 & 0.1224 & 0.1171 & 23.66\% & 20.43\% \\
$(.3,.6]$ & 162 & 0.1279 & 0.1189 & 0.62\% & 1.23\% \\
$(.6,1]$ & 153 & 0.0796 & 0.0749 & 0.00\% & 0.00\% \\
$>1$ & 27 & 0.1115 & 0.1006 & 0.00\% & 0.00\% \\
\bottomrule
\end{tabular}
\end{inlinetable}

表\ref{tab:sources}给出八个场景的结果，表\ref{tab:gains}列出四通道短窗扰动增益。配对场景 Bootstrap 分析先在每个源场景内平均两条偏移轨迹，再以随机种子 20261101 对八个等权场景进行 20,000 次有放回重采样。第 2.5--97.5 百分位区间描述固定开发选定工作点下，这些开发场景之间的结果变动。

完整经验偏移子集的恒零校正 MAE 为 0.548180 DN，恒 $+0.5$ 校正 MAE 为 0.729488 DN。表\ref{tab:buckets}按符号及绝对目标幅度区间统计，最小非零偏移区间 $(0,0.1]\DN$ 含 13 个相互依赖的分量并据此计算误差增大比例 $W$，单独的零偏移条件组衡量精确零偏移行为。

图\ref{fig:z63_empirical}的连续序列以开发集 RGB 辐射代理生成，场景中位亮度为 2 DN。四通道偏移取自 Z63 暗场记录 DSC4194 与 DSC4200 的原生 DN 近邻状态转移模型，每个通道另有序列恒定的有符号分量。八个场景各含 120 帧，评价中心帧为 6--113。相对 T1，T7 在八个场景均降低 $J$，在七个场景降低 $A$。


\section{七帧滤波对照}
\label{app:filters}
连续序列评测基于八个场景的 120 帧序列，包含逐 CFA 恒定本底并对 108 个有效中心帧进行评分。包含阶跃干预的设计压力序列已在正文第 V-D 节与图\ref{fig:designed_regimes}中单独分析。连续预测与真值记录保存在 \texttt{Z63\_CONTINUOUS\_105K.npz}，16 个短窗经验片段则保存在 \texttt{DEV\_105K\_FULL.json}。

均值与中值使用全部七个 T1 值，对称指数权重在七帧片段内正比于 $\rho^{|i-j|}$。Kalman 平滑器采用测量方差 $r_c$、过程方差 $qr_c$ 的高斯随机游走，初始方差为 $10^6r_c$。跳变模型将每次转移定义为
\begin{equation}
 \eta_{k,c}\sim(1-\pi)\mathcal N(0,qr_c)+\pi\mathcal N(0,(q+v)r_c).
\end{equation}
后验均值枚举全部 $2^6=64$ 种转移指示模式，每个状态均由相同七次观测条件化。连续评分与响应测量均只读取滑动窗中心状态。120 帧配对响应序列共用基础序列的几何、光子和读出噪声抽样，在整数量化前加入扰动。共同扰动为各 CFA 恒定增加 0.3 DN，相对扰动将各通道经验轨迹按第 6--113 帧的均值中心化，再将整段 120 帧的最大绝对幅值归一至 0.3 DN。共同增益取配对输出变化的均值，相对增益对中心化配对输出变化与已知扰动作回归，两者均使用与 $A$、$J$ 相同的 108 个中心帧。

每次留出一个场景，使用另外七个场景的零偏移 T1 误差估计 $r_c$。固定网格包含七个指数系数、八个随机游走方差比和 36 组跳变参数，完整列表见 \texttt{FILTER\_PROTOCOL.json}。表\ref{tab:filters_raw}给出响应校准前的结果。为使有限扰动下的响应接近 T7，各类滤波器均增加相同的逐 CFA 校准
\begin{equation}
 \begin{aligned}
 z_{t,c}&=\beta_c[(1-\alpha_c)p_{t,c}+\alpha_cF_c(p)_t],\\
 &0\le\alpha_c\le1,\qquad0\le\beta_c\le1.5.
\end{aligned}
\end{equation}
两个系数由另外七个场景的共同与相对增益目标求解，最大逐 CFA 平均增益差不超过 0.05 的候选按连续 $J$、$A$ 依次排序，候选集为空时选取响应差最小者。五类滤波器在全部八折中均满足训练场景平均增益容限，随后固定系数与参数并应用于留出场景。在留出场景预测中，五类滤波器在扰动测试下的逐 CFA 增益与 T7 的最大差为 0.0353，单个场景与通道组合的最大差为 0.1729。随附数据完整记录了逐场景增益、校准系数、得分与候选工作点。

\begin{table*}[!t]
\centering
\caption{逐 CFA 响应校准前的五类七帧滤波对照。参数由另外七个场景选择，均值与中值采用固定窗口。}\label{tab:filters_raw}
\footnotesize
\begin{tabular}{lrrrr}
\toprule
方法 & $A$ (DN) & $J$ (DN) & $g_r$: min--max & $g_c$: min--max \\
\midrule
T1 & 0.1328 & 0.0605 & 0.718--0.865 & 0.461--0.741 \\
T1 + 均值 & 0.1640 & 0.1154 & 0.082--0.125 & 0.461--0.741 \\
T1 + 中值 & 0.1677 & 0.1133 & 0.045--0.104 & 0.462--0.742 \\
T1 + 指数 & 0.1344 & 0.0661 & 0.582--0.714 & 0.461--0.741 \\
T1 + Kalman & 0.1399 & 0.0775 & 0.461--0.576 & 0.461--0.741 \\
T1 + 跳变模型 & 0.1322 & 0.0531 & 0.696--0.852 & 0.461--0.741 \\
T7 & 0.1270 & 0.0480 & 0.683--0.831 & 0.462--0.743 \\
\bottomrule
\end{tabular}
\end{table*}

\figwide{filter_frontier_zh}{.98}{T1、T7 与五类七帧后处理的开发集工作点。左：绝对误差与相邻变化误差，灰点为未校准网格，彩色标记为响应校准后的选择，颜色表示平均相对增益。右：实测流式耗时与相邻变化误差。T1 缓存逐帧预测，T7 跨窗口缓存空间表征。}{fig:filter_frontier}

相对校准跳变平滑器，T7 在六个场景降低 $J$、七个场景降低 $A$。$J_{T7}-J_{jump}$ 与 $A_{T7}-A_{jump}$ 的配对场景 Bootstrap 95\% 区间分别为 $[-0.010941,-0.002215]\DN$ 和 $[-0.008743,-0.003141]\DN$，由八个开发场景的 20,000 次配对重采样计算，各折预测保持固定。其余四类校准滤波器在全部八个场景的 $J$ 高于 T7。

计时平台为 RTX 3090 Ti、Ryzen 9 9900X、PyTorch 2.6/CUDA 12.4，采用 FP32 与 $320\times480$ Bayer 帧，各分支输出均与原始前向逐位一致。每种模式预热十次，再以交替分组方式计时一百次。T1 缓存逐帧预测，T7 保留六帧空间表征并编码新帧。GPU 计时覆盖模型分支，解码、输入适配和 ISP 单独处理。后处理 CPU 计时采用单线程并预计算协方差矩阵，报告的计算耗时为 GPU 模型与 CPU 滤波实测之和，输出传输与三帧采集前视另计。完整计时样本、显存峰值与独立滤波耗时存于 \texttt{NEURAL\_RUNTIME\_105K.json} 和 \texttt{FILTER\_COMPARISON\_105K.json}。

\section{实拍评价与全分辨率推断}
颜色指标在固定亮度匹配与显影流程下计算相对于参考图像的 CAM16-UCS $a'b'$ 距离均值，主比较采用 S5 II、OM-1 与 S9 的 +6 EV 图像。线性 Bayer 单元值经面积缩小至 768 像素显示宽度后执行固定白平衡、颜色变换和 sRGB 传递函数。平移配准与统一亮度缩放由未校正图像估计并共同应用于校正前后图像。三台相机（S5 II、OM-1 与 S9）均采用相同的预训练模型与显影流程完成评测，其中 S9 采用电子快门与 ISO 100。

一次全分辨率推断对七帧 $6060\times3410$ Bayer 同时输出 T1 和 T7 耗时 5.946328 秒，张量占用的 GPU 显存峰值为 1,753,204,736 字节（1.63 GiB），显存分配器预留 2,470,445,056 字节。使用全部七帧输出时摊销为 0.849475 秒／帧，而步长为一、逐窗独立重算并输出中心帧时为 5.946328 秒／输出帧（计时覆盖首次前向计算区间，解码与 ISP 位于计时区间之外）。

\section{复现材料}
配套归档包含双语源文件与 PDF、105k 评价的冻结预测和真值、PTC 矩阵、曝光阶梯指标，以及可执行滤波分析和图表生成代码，SHA-256 清单连接数值输入与输出。\texttt{neural/} 提供推断实现，\texttt{weights/} 提供预训练 37524、主时域工作点 105k 及图\ref{fig:z63_inference_archive}所用 46k 模型，身份见 \texttt{neural/PROVENANCE.json}。数组接口接收以 DN 表示、减去参考黑电平的带符号 Bayer 打包数据。神经网络响应推断流程记录于 analysis/，随附的预测数组可直接用于滤波基准的复算验证。
